\section{Design and methodology}

This chapter describes the design methodology used to evaluate alternative approaches and select the most suitable one.

\subsection{Evaluation criteria}

The following criteria are used to compare alternatives:
\begin{itemize}
  \item performance --- execution speed under typical workloads;
  \item accuracy --- correctness of produced results;
  \item scalability --- ability to handle increasing data volumes;
  \item maintainability --- ease of extending and modifying the system.
\end{itemize}

% === EXAMPLE: Footnote ===
The evaluation methodology follows the Analytic Hierarchy Process (AHP)\footnote{The AHP was introduced by Thomas L. Saaty in the 1970s and is widely used for multi-criteria decision making.}, which provides a structured framework for comparing alternatives across multiple criteria.

\subsection{Alternatives}

Three alternative design approaches are considered:
\begin{enumerate}
  \item Monolithic architecture --- a single, integrated system;
  \item Microservices architecture --- distributed, independently deployable services;
  \item Hybrid architecture --- a combination of monolithic core with modular extensions.
\end{enumerate}

\subsection{Comparison matrix}

Table~\ref{tab:alternatives} presents a comparison of the three alternatives.

% === EXAMPLE: Complex table with multirow header ===
\begin{table}[H]
    \caption{Comparison of design alternatives} \label{tab:alternatives}
    \begin{tabular}{|l|c|c|c|c|}
    \hline \textbf{Criterion} & \textbf{Weight} & \textbf{Monolithic} & \textbf{Microservices} & \textbf{Hybrid} \\
    \hline Performance  & 0.30 & 8 & 7 & 9 \\
    \hline Accuracy     & 0.25 & 9 & 8 & 9 \\
    \hline Scalability  & 0.25 & 5 & 9 & 8 \\
    \hline Maintainability & 0.20 & 6 & 8 & 7 \\
    \hline \textbf{Weighted score} & --- & \textbf{7.15} & \textbf{7.95} & \textbf{8.30} \\
    \hline
    \end{tabular}
\end{table}

\subsection{Selected approach}

Based on the weighted comparison, the \textbf{hybrid architecture} achieves the highest overall score of 8.30.
This approach combines the performance benefits of a monolithic core with the scalability advantages of modular extensions.

The effectiveness $E$ of each alternative is calculated using Equation~\eqref{eq:basic} from the literature review (Section~2.2):

% === EXAMPLE: Cross-referencing equations from other chapters ===
\begin{equation} \label{eq:weighted}
    E = \sum_{j=1}^{m} w_j \cdot s_j
\end{equation}

where $w_j$ is the weight of criterion $j$ and $s_j$ is the score assigned to that criterion.

% === EXAMPLE: TikZ Gantt-style timeline ===
The implementation timeline is shown in Figure~\ref{fig:timeline}.

\begin{figure}[H]
    \centering
    \begin{tikzpicture}[>=Stealth]
        % Timeline axis
        \draw[->] (0,0) -- (13,0) node[right] {Weeks};
        \foreach \x in {1,...,12} {
            \draw (\x,-0.1) -- (\x,0.1) node[below=3pt] {\small \x};
        }

        % Task bars
        \fill[blue!40] (1,0.5) rectangle (3,1) node[midway, font=\small] {Analysis};
        \fill[green!40] (3,0.5) rectangle (6,1) node[midway, font=\small] {Design};
        \fill[orange!40] (5,1.2) rectangle (9,1.7) node[midway, font=\small] {Implementation};
        \fill[red!30] (8,0.5) rectangle (11,1) node[midway, font=\small] {Testing};
        \fill[purple!30] (10,1.2) rectangle (12,1.7) node[midway, font=\small] {Writing};
    \end{tikzpicture}
    \caption{Project timeline (Gantt chart)} \label{fig:timeline}
\end{figure}

% === EXAMPLE: Description list ===
\textbf{Advantages of the hybrid approach:}
\begin{description}
    \item[Performance] The monolithic core avoids network overhead for critical operations.
    \item[Scalability] Extension modules can be independently scaled based on demand.
    \item[Maintainability] Clear boundaries between core and extensions simplify updates.
    \item[Deployment] Docker containers enable consistent deployment across environments.
\end{description}

\clearpage
